\section{Exercícios N1--N4}

\begin{frame}{N1 — leitura da planta e requisitos mínimos}
\small
\begin{enumerate}
\item Para uma sala de 18 m$^2$ e perímetro de 17 m, determine a previsão mínima de carga de iluminação e o número mínimo de TUG.
\item Explique por que a potência de previsão de iluminação não determina diretamente a quantidade de luminárias.
\item Em uma planta residencial, indique os locais adequados para QD, interruptores e tomadas considerando portas, circulação e mobiliário.
\item Diferencie TUG, TUE, circuito dedicado, PE, DR/DDR e DPS.
\end{enumerate}
\end{frame}

\begin{frame}{N2 — dimensionamento de circuitos}
\small
\begin{enumerate}
\item Dimensione um circuito de TUG em 220 V, 1.900 VA e 24 m, verificando corrente, proteção, seção inicial e queda de tensão.
\item Analise um chuveiro de 5,5 kW em 220 V a 28 m do QD; verifique a seção por ampacidade e queda de tensão.
\item Proponha a divisão de circuitos para uma residência com sala, cozinha, área de serviço, dois quartos, banheiro e escritório.
\item Explique quais desses circuitos exigem proteção diferencial-residual de alta sensibilidade e por quê.
\end{enumerate}
\end{frame}

\begin{frame}{N3 — quadro, balanceamento e entrada Equatorial}
\small
\begin{enumerate}
\item Distribua 11 circuitos monofásicos de potências distintas entre A, B e C em uma rede 380/220 V, minimizando o desbalanceamento da carga instalada.
\item Para uma unidade residencial no Maranhão com 21,5 kW instalados, determine o enquadramento de fornecimento usando a NT.00001.EQTL Rev.09 e identifique quais dados ainda precisam ser verificados no padrão de entrada.
\item Mostre por que o condutor mínimo indicado pela concessionária para o padrão não pode ser copiado automaticamente para o alimentador interno.
\item Desenhe um unifilar com DG, DPS, barramentos N/PE, circuitos com DDR/RCBO e identificação coerente com o quadro de cargas.
\end{enumerate}
\end{frame}

\begin{frame}{N4 — desafio de projeto residencial completo}
\small
Receba uma planta arquitetônica residencial cotada e, sem alterar sua geometria, desenvolva o projeto elétrico completo:
\begin{enumerate}
\item calcule os pontos mínimos de iluminação e TUG e justifique pontos adicionais;
\item posicione TUE e cargas específicas conforme uso real;
\item divida e identifique os circuitos na planta;
\item dimensione condutores, eletrodutos e proteções;
\item verifique queda de tensão, proteção diferencial, DPS, PE e equipotencialização;
\item balanceie as fases e enquadre a entrada conforme Equatorial Maranhão;
\item gere quadro de cargas, unifilar, memorial de cálculo e lista de materiais;
\item confira se planta, quadro e unifilar apresentam exatamente os mesmos circuitos.
\end{enumerate}
\end{frame}

%=====================================================================

